\documentclass[../../../main]{subfiles}

\ifSubfilesClassLoaded{%
    \zexternaldocument*{../../1.Category-Theory/exported/aux/main}%
    \zexternaldocument*{../../2.Homological-Algebra/exported/aux/main}%
}{}

\begin{document}
\section{Presheaves} \label{sec:presheaves}

\begin{definition}[Presheaf] \label{def:presheaf}
  Let \(X\) be a topological space and let \(\Open(X)\) be its \hyperref[ex:open-set-category]{posetal category of open sets}.
  Let \(\category{C}\) be a \hyperref[def:category]{category}.
  The \hyperref[def:category]{category} of \(\category{C}\)-valued \textbf{presheaves} over \(X\) is the \hyperref[def:functor-category]{functor category}
  \begin{equation*}
    \PSh_{X}(\category{C}) \coloneq \Fun{\Open(X)^{\op}}{\category{C}}
  \end{equation*}
  The elements of this category are \hyperref[def:functor]{functors} \(\mathcal{F}: \Open(X)^{\op} \to \category{C}\) and morphisms are \hyperref[def:natural-transformation]{natural transformations} \(\phi: \mathcal{F} \to \mathcal{G}\).
\end{definition}

Less abstractly, a \(\category{C}\)-valued \hyperref[def:presheaf]{presheaf} \(\mathcal{F}\) is the collection of data:
\begin{itemize}
  \item For each open set \(U \subseteq X\) an object \(\mathcal{F}(U) \in \category{C}\) called \textit{sections of \(\mathcal{F}\) over \(U\)}.
  \item For every inclusion of open sets \(V \subseteq U\) a morphism \(\res_{U, V}: \mathcal{F}(U) \to \mathcal{G}(U)\) in \(\category{C}\) called \textit{restriction} such that
        \begin{itemize}
          \item For every open set \(U \subseteq X\), \(\res_{U, U} = \id_{U}\).
          \item For every inclusion of open sets \(W \subseteq V \subseteq U\), \(\res_{U, V} \circ \res_{V, W} = \res_{U, W}\).
        \end{itemize}
\end{itemize}
A morphism of \hyperref[def:presheaf]{presheaves} \(\phi: \mathcal{F} \to \mathcal{G}\) is a collection of morphisms
\begin{equation*}
  \{\phi(U): \mathcal{F}(U) \to \mathcal{G}(U)\}_{U \in \Open(X)}
\end{equation*}
in \(\category{C}\) making the following diagram commute
\begin{center}
  \begin{tikzcd}
    \mathcal{F}(U) \arrow[r, "\phi(U)"] \arrow[d, "\res_{U, V}"'] & \mathcal{G}(U) \arrow[d, "\res_{U, V}"] \\
    \mathcal{F}(V) \arrow[r, "\phi(V)"'] & \mathcal{G}(V)
  \end{tikzcd}
\end{center}

\begin{notation}
  When the context is clear, we will drop the subscript of \(\PSh_{X}(\category{C})\) and denote the category of \(\category{C}\)-valued \hyperref[def:presheaf]{presheaves} over \(X\) as \(\PSh(\category{C})\).
\end{notation}

Since \(\PSh_{X}(\category{C})\) is a \hyperref[def:functor-category]{functor category}, if \(\category{C}\) has \hyperref[def:categorical-limit]{limits} or \hyperref[def:categorical-colimit]{colimits} of a certain shape then \(\PSh_{X}(\category{C})\) also has them and they are computed pointwise.
This means that \hyperref[def:categorical-limit]{limits} and \hyperref[def:categorical-colimit]{colimits} can be computed \textit{open set by open set}.

\begin{proposition} \label{prop:limits-colimits-presheaf}
  Let \(X\) be a topological space.
  Let \(\mathcal{I} \to \PSh_{X}(\category{C}), i \mapsto \mathcal{F}_{i}\) be a \hyperref[def:diagram]{diagram}.
  If \(\category{C}\) has \hyperref[def:categorical-limit]{limits} (resp. \hyperref[def:categorical-colimit]{colimits}) of shape \(\mathcal{I}\) then \(\lim_{i} \mathcal{F}_{i}\) (resp. \(\colim_{i} \mathcal{F}_{i}\)) exists in \(\PSh_{X}(\category{C})\).
  Moreover for every open set \(U \subseteq X\)
  \begin{align*}
    \left( \lim_{i} \mathcal{F}_{i} \right)(U) = \lim_{i} \left( \mathcal{F}_{i}(U) \right) \qquad \left( \colim_{i} \mathcal{F}_{i} \right)(U) = \colim_{i} \left( \mathcal{F}_{i}(U) \right)
  \end{align*}
\end{proposition}
\begin{proof}
  Particular case of \autoref{prop:limits-in-functor-category}.
\end{proof}


In particular if \(\category{C}\) is an \hyperref[def:abelian-category]{abelian category} then \hyperref[def:categorical-kernel]{kernels} and \hyperref[def:categorical-cokernel]{cokernels} exist in \(\PSh_{X}(\category{C})\) and are computed \textit{open set by open set}, making \(\PSh_{X}(\category{C})\) an \hyperref[def:abelian-category]{abelian category}.

\begin{example} \label{ex:one-point-presheaf}
  Consider \(X = \{\ast\}\) the one-point topological space.
  The category \(\Open(X)\) has
  \begin{align*}
    \operatorname{Ob}(\Open(X)) = \{\varnothing, X \} \qquad \Hom(\Open(X)) = \{\varnothing \hookrightarrow X\}
  \end{align*}
  Therefore every \(\mathcal{F} \in \PSh_{X}(\category{C})\) consists of two objects \(\mathcal{F}(\varnothing), \mathcal{F}(X) \in \category{C}\) and a restriction morphism \(\rho: \mathcal{F}(X) \to \mathcal{F}(\varnothing)\).

  This makes \(\PSh_{\{\ast\}}(\category{C}) \cong \mathsf{Arr}(\category{C})\).
\end{example}


\subsection{Presheaves and continuous maps} \label{ssec:presheaves-and-continuous-maps}
\begin{definition}[Direct image presheaf] \label{def:direct-image-presheaf}
  Let \(f: X \to Y\) be a continuous map between two topological spaces \(X, Y\).
  Let \(\mathcal{F}\) be a \(\category{C}\)-valued \hyperref[def:presheaf]{presheaf} over \(X\).

  The \textbf{direct image presheaf} of \(\mathcal{F}\) via \(f\), \(f_{\ast}\mathcal{F} \in \PSh_{Y}(\category{C})\), is defined as
  \begin{equation*}
    \forall V \in \Open(Y) \quad f_{\ast}\mathcal{F}(V) = \mathcal{F}(f^{-1}V).
  \end{equation*}
  Given open subsets \(W \subseteq V \subseteq Y\), the restriction maps \(\res_{V, W}\) comes directly from the diagram
  \begin{center}
    \begin{tikzcd}
      f_{\ast}\mathcal{F}(V) \arrow[r, equal] \arrow[d, dotted] & \mathcal{F}(f^{-1} V) \arrow[d, "\res_{f^{-1}V, f^{-1}W}"] \\
      f_{\ast}\mathcal{F}(W) \arrow[r, equal] & \mathcal{F}(f^{-1}W)
    \end{tikzcd}
  \end{center}

  Notice that this construction defines a \hyperref[def:functor]{functor} \(f_{\ast}: \PSh_{X}(\category{C}) \to \PSh_{Y}(\category{C})\).
\end{definition}

\begin{example}[Global section] \label{ex:global-section-presheaf}
  Let \(X\) be a topological space.
  Consider the one-point topological space \(\{\ast\}\) and the unique continuous map \(f: X \to \{\ast\}\).
  Let \(\mathcal{F}\) be a \(\category{C}\)-valued \hyperref[def:presheaf]{presheaf} over \(X\).

  Consider the \hyperref[def:direct-image-presheaf]{direct image presheaf} of \(\mathcal{F}\) via \(f\).
  Since this is a \(\category{C}\)-valued \hyperref[def:presheaf]{presheaf} over \(\{\ast\}\) by \autoref{ex:one-point-presheaf} we have that
  \begin{align*}
    f_{\ast}\mathcal{F} \cong \left( \mathcal{F}(\{\ast\}) \to \mathcal{F}(\varnothing) \right)
  \end{align*}
  Using the definition of the \hyperref[def:direct-image-presheaf]{direct image presheaf} we have
  \begin{align*}
    f_{\ast}\mathcal{F}(\varnothing) &= \mathcal{F}(f^{-1}(\varnothing)) = \mathcal{F}(\varnothing) \\
    f_{\ast}\mathcal{F}(\{\ast\}) &= \mathcal{F}(f^{-1}(\{\ast\})) = \mathcal{F}(X)
  \end{align*}

  Therefore we define the \textbf{global sections} of \(\mathcal{F}\) as
  \begin{align*}
    \Gamma(X, \mathcal{F}) \coloneq f_{\ast}\mathcal{F} \cong \mathcal{F}(X) \to \mathcal{F}(\varnothing)
  \end{align*}

  This defines a \hyperref[def:functor]{functor} \(\Gamma(X, -): \PSh_{X}(\category{C}) \to \mathsf{Arr}(\category{C})\).

  In particular, if \(\mathcal{F}(\varnothing)\) is a \hyperref[def:initial-terminal-object]{terminal object} in \(\category{C}\) we can identify \(\mathcal{F}(X) \to \mathcal{F}(\varnothing)\) with \(\mathcal{F}(X)\), getting a \hyperref[def:functor]{functor} \(\Gamma(X, -): \PSh_{X}(\category{C}) \to \category{C}\).
\end{example}

\begin{example}[Skyscraper presheaf]\label{ex:skyscraper-presheaf}
  Let \(X\) be a topological space and \(p \in X\).
  Consider the one-point topological space \(\{\ast\}\) and the continuous map \(\iota_{p}: \{\ast\} \to X, \ast \mapsto p\).
  Consider a morphism \(\rho: S \to T\) in \(\category{C}\), which by \autoref{ex:one-point-presheaf} corresponds to a \hyperref[def:presheaf]{presheaf} \(\mathcal{R} \in \PSh_{\{\ast\}}(\category{C})\).

  Consider the \hyperref[def:direct-image-presheaf]{direct image presheaf} of \(\mathcal{R}\) via \(\iota_{p}\):
  \begin{align*}
    \forall U \in \Open(X) \quad \iota_{p, \ast}\mathcal{R}(U) = \mathcal{R}(\iota_{p}^{-1}(U)) = \begin{cases}
      \mathcal{R}(\{\ast\}) & p \in U \\
      \mathcal{R}(\varnothing) & p \notin U
    \end{cases} = \begin{cases}
      S & p \in U \\
      T & p \notin U
    \end{cases}
  \end{align*}
  Note that if \(V \subseteq U\) the restriction map is given by
  \begin{align*}
    \res_{U, V} = \begin{cases}
      \id_{S} & p \in V \\
      \id_{T} & p \notin U \\
      \rho & p \in U \land p \notin V
    \end{cases}
  \end{align*}

  We define the \textbf{skyscraper presheaf} supported at \(\rho\), \(\iota_{p, \ast}\rho \in \PSh_{X}(\category{C})\), as
  \begin{equation}
    \iota_{p, \ast} \underline{\rho}_{\operatorname{pre}} \coloneq \iota_{p, \ast} \mathcal{R}
  \end{equation}
  This defines a \hyperref[def:functor]{functor} \(\iota_{p, \ast}: \mathsf{Arr}(\category{C}) \to \PSh_{X}(\category{C})\).

  In particular if \(T = \mathcal{R}(\varnothing)\) is a \hyperref[def:initial-terminal-object]{terminal object} in \(\category{C}\) we can identify \(\rho: S \to T\) with \(S\), getting a \hyperref[def:functor]{functor} \(\iota_{p, \ast}: \category{C} \to \PSh_{X}(\category{C})\).
  In this situation we write \(\iota_{p, \ast} \underline{S}_{\operatorname{pre}}\).
\end{example}


For sufficiently nice \hyperref[def:category]{categories} (those with \hyperref[def:direct-limit]{direct limits}), we have a \textit{dual} construction.
\begin{definition}[Inverse image presheaf] \label{def:inverse-image-presheaf}
  Let \(f: X \to Y\) be a continuous map between two topological spaces \(X, Y\).
  Let \(\category{C}\) be a \hyperref[def:category]{category} with \hyperref[def:direct-limit]{direct limits}.
  Let \(\mathcal{G}\) be a \(\category{C}\)-valued \hyperref[def:presheaf]{presheaf} over \(Y\).

  The \textbf{inverse image presheaf} of \(\mathcal{G}\), \(f_{\operatorname{pre}}^{-1}\mathcal{G} \in \PSh_{X}(\category{C})\), of \(\mathcal{G}\) via \(f\) is the \hyperref[def:presheaf]{presheaf}
  \begin{equation*}
    \forall U \in \Open(X) \quad \left( f_{\operatorname{pre}}^{-1}\mathcal{G} \right)(U) \coloneq \varinjlim_{\substack{V \supseteq f(U) \\ V \in \Open(Y) }} \mathcal{G}(V)
  \end{equation*}

  Notice that this defines a \hyperref[def:functor]{functor} \(f_{\operatorname{pre}}^{-1}: \PSh_{Y}(\category{C}) \to \PSh_{X}(\category{C})\).
\end{definition}

\begin{remark} \label{remark:inverse-image-presheaf-set}
  For \(\category{C} = \Set\) we can describe \(f^{-1}\mathcal{G}(U)\) as
  \begin{equation*}
    \left( f_{\operatorname{pre}}^{-1} \mathcal{G} \right)(U) = \{ (s, V) \mid f(U) \subseteq V \in \Open(Y) \land s \in \mathcal{G}(V) \}
  \end{equation*}
  modulo the relation
  \begin{equation*}
    (s, V) \sim (s', V') \iff \exists W \in \Open(Y) \mid W \subseteq V \cap V' \land f(U) \subseteq W \land s|_{W} = s'|_{W}
  \end{equation*}
\end{remark}

In the following examples, we assume that the \hyperref[def:category]{category} \(\category{C}\) has \hyperref[def:direct-limit]{direct limits}.
\begin{example}[Constant presheaf] \label{ex:constant-presheaf}
  Let \(X\) be a topological space.
  Consider the one-point topological space \(\{\ast\}\).
  Consider a morphism \(\rho: S \to T\) in \(\category{C}\).
  Note that by \autoref{ex:one-point-presheaf} this corresponds to a \hyperref[def:presheaf]{presheaf} \(\mathcal{R} \in \PSh_{\{\ast\}}(\category{C})\).

  Consider the \hyperref[def:inverse-image-presheaf]{inverse image presheaf} of \(\mathcal{R}\) via the map \(f: X \to \{\ast\}\).
  \begin{itemize}
    \item If \(U \in \Open(X)\) is not the empty set we have that
          \begin{align*}
            (f_{\operatorname{pre}}^{-1}\mathcal{R})(U) = \varinjlim_{\substack{V \supseteq f(U) \\ V \in \Open(\{\ast\})}} \mathcal{R}(V) = \varinjlim_{\substack{V \ni \ast \\ V \in \Open(\{\ast\})}} \mathcal{R}(V) = \mathcal{R}(\{\ast\}) = S
          \end{align*}
    \item If \(U = \varnothing\) then we take the \hyperref[def:categorical-limit]{limit} over the diagram \(S \xrightarrow{\rho} T\), getting
          \begin{equation*}
            (f_{\operatorname{pre}}^{-1}\mathcal{R})(\varnothing) = T
          \end{equation*}
  \end{itemize}
  Therefore
  \begin{equation*}
    \forall U \in \Open(X) \quad \left( f_{\operatorname{pre}}^{-1}\mathcal{R} \right)(U) = \begin{cases}
      S & U \neq \varnothing \\
      T & U = \varnothing
    \end{cases}
  \end{equation*}

  Note that for \(V \subseteq U\) open subsets of \(X\) the restriction maps are given by
  \begin{align*}
    \res_{U, V} = \begin{cases}
      \id_{S} & V \neq \varnothing \\
      \id_{T} & U = \varnothing \\
      \rho & V = \varnothing \land U \neq \varnothing
    \end{cases}
  \end{align*}


  Therefore we can define the \textbf{constant presheaf} supported as \(\rho\) as
  \begin{equation*}
    \underline{\rho}_{\operatorname{pre}}(U) \coloneq  f_{\operatorname{pre}}^{-1}\mathcal{R}
  \end{equation*}
  This defines a \hyperref[def:functor]{functor} \(\underline{\cdot}_{\operatorname{pre}}: \mathsf{Arr}(\category{C}) \to \PSh_{X}(\category{C})\)

  Note that if \(T = \mathcal{R}(\varnothing)\) is a \hyperref[def:initial-terminal-object]{terminal object} in \(\category{C}\) we can identify \(\rho: S \to T\) with \(S\), getting a \hyperref[def:functor]{functor} \(\underline{\cdot}_{\operatorname{pre}}: \category{C} \to \PSh_{X}(\category{C})\).
  In this situation we write \(\underline{S}_{\operatorname{pre}}\).
\end{example}

\begin{example}[Stalk]\label{ex:stalk-presheaf}
  Let \(X\) be a topological space and \(p \in X\).
  Consider the one-point topological space \(\{\ast\}\) and the inclusion \(f: \{\ast\} \hookrightarrow X, \ast \mapsto p\).
  Let \(\mathcal{F}\) be a \(\category{C}\)-valued \hyperref[def:presheaf]{presheaf} over \(X\).

  Consider the \hyperref[def:inverse-image-presheaf]{inverse image presheaf} of \(\mathcal{F}\) via \(f\).
  Since this is a \(\category{C}\)-valued \hyperref[def:presheaf]{presheaf} over \(\{\ast\}\) by \autoref{ex:one-point-presheaf}
  \begin{align*}
    f_{\operatorname{pre}}^{-1} \mathcal{F} \cong \left( f_{\operatorname{pre}}^{-1}\mathcal{F}(\{\ast\}) \to f_{\operatorname{pre}}^{-1}\mathcal{F}(\varnothing) \right)
  \end{align*}
  Using the definition of the \hyperref[def:inverse-image-presheaf]{inverse image presheaf} we have that
  \begin{align*}
    f_{\operatorname{pre}}^{-1}\mathcal{F}(\{\ast\}) &= \varinjlim_{\substack{U \supseteq f(\{\ast\}) \\ U \in \Open(X)}} \mathcal{F}(U) = \varinjlim_{\substack{U \ni p \\ U \in \Open(X)}} \mathcal{F}(U) \\
    f_{\operatorname{pre}}^{-1}\mathcal{F}(\varnothing) &= \varinjlim_{\substack{U \supseteq f(\varnothing) \\ U \in \Open(X)}} \mathcal{F}(U) = \varinjlim_{\substack{U \in \Open(X)}} \mathcal{F}(U) = \mathcal{F}(\varnothing)
  \end{align*}
  Therefore
  \begin{align*}
    f_{\operatorname{pre}}^{-1} \mathcal{F} \cong \left( \varinjlim_{\substack{U \ni p \\ U \in \Open(X)}} \mathcal{F}(U) \to \mathcal{F}(\varnothing) \right)
  \end{align*}

  We define the \textbf{stalk} of \(\mathcal{F}\) at \(p\) as
  \begin{equation*}
    \mathcal{F}_{p} \coloneq f_{\operatorname{pre}}^{-1}\mathcal{F}
  \end{equation*}
  This defines a \hyperref[def:functor]{functor} \(-_{p}: \PSh_{X}(\category{C}) \to \mathsf{Arr}(\category{C})\).

  In particular if \(\mathcal{F}(\varnothing)\) is a \hyperref[def:initial-terminal-object]{terminal object} in \(\category{C}\) we can identify
  \begin{equation*}
    \left(\varinjlim_{\substack{U \ni p \\ U \in \Open(X)}} \mathcal{F}(U) \to \mathcal{F}(\varnothing) \right) \cong \varinjlim_{\substack{U \ni p \\ U \in \Open(X)}} \mathcal{F}(U)
  \end{equation*}
  getting a \hyperref[def:functor]{functor} \(-_{p}: \PSh_{X}(\category{C}) \to \category{C}\).
\end{example}

\begin{lemma}
  Let \(f: X \to Y\) be a continuous map between two topological space \(X, Y\) and \(p \in X\).
  Let \(\category{C}\) be a \hyperref[def:category]{category} with \hyperref[def:direct-limit]{direct limits}.
  Let \(\mathcal{G}\) be a \(\category{C}\)-valued \hyperref[def:presheaf]{presheaf} over \(Y\).
  Then \((f_{\operatorname{pre}}^{-1} \mathcal{G})_{p} = \mathcal{G}_{f(p)}\).
\end{lemma}
\begin{proof}
  \begin{align*}
    (f_{\operatorname{pre}}^{-1} \mathcal{G})_{p} &= \varinjlim_{\substack{U \ni p \\ U \in \Open(X)}} f_{\operatorname{pre}}^{-1}\mathcal{G}(U) \to \varinjlim_{U \in \Open(X)} f_{\operatorname{pre}}^{-1}\mathcal{G}(U) \\
                                        &= \varinjlim_{\substack{U \ni p \\ U \in \Open(X)}} \left( \varinjlim_{\substack{V \supseteq f(U) \\ V \in \Open(Y)}} \mathcal{G}(V) \right) \to \varinjlim_{U \in \Open(X)} \left( \varinjlim_{\substack{V \supseteq f(U) \\ V \in \Open(Y)}} \mathcal{G}(V) \right) \\
                                        &= \varinjlim_{\substack{V \ni f(p) \\ V \in \Open(Y)}} \mathcal{G}(V) \to \varinjlim_{V \in \Open(Y)} \mathcal{G}(V) \\
                                        &= \mathcal{G}_{f(p)}
  \end{align*}
\end{proof}


\begin{theorem} \label{thm:direct-inverse-image-adjunction-presheaf}
  Let \(f: X \to Y\) be a continuous function between two topological spaces \(X, Y\).
  Let \(\category{C}\) be a \hyperref[def:category]{category} with \hyperref[def:direct-limit]{direct limits}.
  Then there exists an \hyperref[def:adjoint-functors]{adjunction}
  \begin{center}
    \begin{tikzcd}
      f_{\operatorname{pre}}^{-1} \dashv f_{\ast}: \PSh_{X}(\category{C}) \arrow[r, shift right, "f_{\ast}"'] &  \PSh_{Y}(\category{C}) \arrow[l, shift right, "f_{\operatorname{pre}}^{-1}"']
    \end{tikzcd}
  \end{center}
\end{theorem}
\begin{proof}
  We want to construct \hyperref[def:natural-transformation]{natural transformations}
  \begin{align*}
    \eta: \id_{\PSh_{Y}(\category{C})} &\implies f_{\ast} \circ f_{\operatorname{pre}}^{-1} \\
    \varepsilon: f_{\operatorname{pre}}^{-1} \circ f_{\ast} &\implies \id_{\PSh_{X}(\category{C})}
  \end{align*}
  making for all \(\mathcal{F} \in \PSh_{X}(\category{C})\) and \(\mathcal{G} \in \PSh_{Y}(\category{C})\) the following diagrams commute
  \begin{center}
    \begin{tikzcd}
      f_{\operatorname{pre}}^{-1}(\mathcal{G}) \arrow[r, "\id_{\mathcal{G}}"] \arrow[d, "f_{\operatorname{pre}}^{-1}(\eta_{\mathcal{G}})"'] & f_{\operatorname{pre}}^{-1}(\mathcal{G}) \\
      \left( f_{\operatorname{pre}}^{-1} \circ f_{\ast} \circ f_{\operatorname{pre}}^{-1} \right)(\mathcal{G}) \arrow[ur, "\varepsilon_{f_{\operatorname{pre}}^{-1}(\mathcal{G})}"']
    \end{tikzcd}
    \begin{tikzcd}
      f_{\ast}(\mathcal{G}) \arrow[r, "\id_{\mathcal{F}}"] \arrow[d, "\eta_{f_{\ast}(\mathcal{F})}"'] & f_{\ast}(\mathcal{G}) \\
      \left( f_{\ast} \circ f_{\operatorname{pre}}^{-1} \circ f_{\ast} \right)(\mathcal{F}) \arrow[ur, "f_{\ast}(\varepsilon_{\mathcal{F}})"']
    \end{tikzcd}
  \end{center}

  Note that, by definition of the \hyperref[def:direct-image-presheaf]{direct image presheaf} and \hyperref[def:inverse-image-presheaf]{inverse image presheaf} we have, for \(U \in \Open(X)\) and \(V \in \Open(Y)\)
  \begin{align*}
    \left( f_{\operatorname{pre}}^{-1} \circ f_{\ast} \right)(\mathcal{F})(U) &= \varinjlim_{\substack{W \supseteq f(U) \\ W \in \Open(Y)}} (f_{\ast}\mathcal{F})(W) = \varinjlim_{\substack{W \supseteq f(U) \\ W \in \Open(Y)}} \mathcal{F}(f^{-1}(W)) \\
    \left( f_{\ast} \circ f_{\operatorname{pre}}^{-1} \right)(\mathcal{G})(V) &= (f_{\operatorname{pre}}^{-1}\mathcal{G})(f^{-1}(V)) = \varinjlim_{\substack{W \supseteq (f \circ f^{-1})(V) \\ W \in \Open(Y)}} \mathcal{G}(W) \end{align*}

  Since we have \((f \circ f^{-1})(V) \subseteq V\), we can define \(\eta_{\mathcal{G}}(V)\) as the canonical injection into the \hyperref[def:direct-limit]{direct limit}
  \begin{equation*}
    \eta_{\mathcal{G}}(V): \mathcal{G}(V) \to \varinjlim_{\substack{W \supseteq (f \circ f^{-1})(V) \\ W \in \Open(Y)}} \mathcal{G}(W)
  \end{equation*}

  On the other hand, since for every \(W \subseteq f(U)\) open we have \(f^{-1}(W) \subseteq (f^{-1} \circ f)(U) \subseteq U\) there exists a restriction morphism
  \begin{equation*}
    \res_{f^{-1}(W), U}: \mathcal{F}(f^{-1}(W)) \to \mathcal{F}(U)
  \end{equation*}
  Using this maps we obtain a morphism
  \begin{equation*}
    \varepsilon_{\mathcal{F}}(U): \varinjlim_{\substack{W \supseteq f(U) \\ W \in \Open(Y)}} \mathcal{F}(f^{-1}(W)) \to \mathcal{F}(U)
  \end{equation*}
\end{proof}

In particular, we get the following particular cases of \autoref{thm:direct-inverse-image-adjunction-presheaf} for a \hyperref[def:category]{category} \(\category{C}\) with \hyperref[def:direct-limit]{direct limits}.
\begin{corollary} \label{cor:constant-global-adjunction-presheaf}
  Let \(X\) be a topological space and \(p \in X\).
  Let \(\mathcal{F}\) be a \(\category{C}\)-valued \hyperref[def:presheaf]{presheaf} over \(X\) and let \(\rho: S \to T\) be a morphism in \(\category{C}\).
  Then
  \begin{align*}
    \Hom_{\PSh_{X}(\category{C})}(\underline{\rho}_{\operatorname{pre}}, \mathcal{F}) \simeq \Hom_{\mathsf{Arr}(\category{C})}(\rho, \Gamma(X, \mathcal{F}))
  \end{align*}

  In particular, if \(\mathcal{F}(\varnothing)\) and \(T\) are \hyperref[def:initial-terminal-object]{terminal objects} in \(\category{C}\) we get
  \begin{align*}
    \Hom_{\PSh_{X}(\category{C})}(\underline{S}_{\operatorname{pre}}, \mathcal{F}) \simeq \Hom_{\category{C}}(S, \Gamma(X, \mathcal{F}))
  \end{align*}
\end{corollary}

\begin{corollary} \label{cor:stalk-skyscraper-adjunction-presheaf}
  Let \(X\) be a topological space and \(p \in X\).
  Let \(\mathcal{F}\) be a \(\category{C}\)-valued \hyperref[def:presheaf]{presheaf} over \(X\) and let \(\rho: S \to T\) be a morphism in \(\category{C}\).
  Then
  \begin{align*}
    \Hom_{\mathsf{Arr}(\category{C})}(\mathcal{F}_{p}, \rho) \simeq \Hom_{\PSh_{X}(\category{C})}(\mathcal{F}, \iota_{p, \ast} \underline{\rho}_{\operatorname{pre}} )
  \end{align*}

  In particular, if \(\mathcal{F}(\varnothing)\) and \(T\) are \hyperref[def:initial-terminal-object]{terminal objects} in \(\category{C}\) we get
  \begin{align*}
    \Hom_{\category{C}}(\mathcal{F}_{p}, S) \simeq \Hom_{\PSh_{X}(\category{C})}(\mathcal{F}, \iota_{p, \ast} \underline{S}_{\pre})
  \end{align*}
\end{corollary}

\end{document}
